\section{Introduction}
\label{sec:intro}

Lean 4 with Mathlib ships strong general-purpose proof automation:
\texttt{aesop} for forward search, \texttt{simp} for rewriting,
\texttt{linarith} and \texttt{polyrith} for linear and polynomial
arithmetic, \texttt{nlinarith} for non-linear arithmetic,
\texttt{positivity} for sign goals, \texttt{measurability} for the
$\sigma$-algebra fragment, \texttt{bound} for monotone composition,
\texttt{gcongr} for generalised congruence
\citep{moura2021lean, mathlib2020community, mathlib2025}.
Applied-mathematics proofs route through these general tactics,
typically by hand decomposition into a chain of \texttt{simp} sets,
\texttt{linarith} hypotheses, and Mathlib lemma selection. None of
the general tactics know about each other. The user is the
orchestrator. There is no domain hammer for applied mathematics.

CoqHammer \citep{czajkakaliszyk2018coqhammer} dispatches Coq goals
to external automated theorem provers (ATPs) such as Vampire, E,
and CVC4, with kernel-checked reconstruction. Sledgehammer
\citep{paulsonblanchette2010threeyears, blanchettesledgehammer2011}
plays the same role in Isabelle/Higher Order Logic (HOL). Lean 4 has no comparable
hammer. This paper does not attempt the fully-general case across
all of Lean. We address the applied-mathematics surface: a
registered library of concentration, anytime-valid, optimal
transport, control-theoretic, and information-theoretic lemmas
plus a fan-out across local Satisfiability Modulo Theories (SMT)
and ATP solvers, packaged as a single deterministic tactic call.

\paragraph{Contribution.}
We present \texttt{pythia}, a Lean 4 tactic library for
applied-mathematics proof automation. The headline tactic
\texttt{pythia!} orchestrates a nine-rung deterministic ladder
under one call. The first five rungs are pure-Lean tactics. They are
\texttt{stat\_simp} (curated probability and ENNReal normal
forms), the \texttt{linarith}/\texttt{nlinarith}/\texttt{polyrith}
chain, \texttt{positivity}, \texttt{aesop} on the registered
\texttt{Pythia} ruleset, and the \texttt{@[stat\_lemma]}
shape-dispatch cascade. The next three rungs are external solvers
with kernel-checked reconstruction. They are \texttt{z3\_check}
(linear real arithmetic via Z3), \texttt{cvc5\_check} (bit-vector
and linear real arithmetic via CVC5), and \texttt{vampire\_check}
or \texttt{e\_check} (first-order logic, FOL, via Vampire or E).
The final rung is \texttt{disprove}, a counterexample finder via
Z3 model extraction. Each rung tries fail-fast with a per-rung
budget, and first-to-close wins. The
verbose variant \texttt{pythia?} reports the closing rung and
per-rung wall-clock timing for diagnostic use. \texttt{pythia}
ships ten registered tactics in total. Six attribute-driven
registries cover lemma routing: \texttt{@[stat\_lemma]},
\texttt{@[stats\_ineq]}, \texttt{@[prob\_simp]},
\texttt{@[stat\_simp]}, \texttt{@[anytime\_valid\_lemma]}, plus
the \texttt{DomainCalculator} typeclass for per-domain numerical
calculators. The proof library spans fifteen domains and carries
thirty-two cross-domain theorems with matched Python sim runners.
Sim runners exercise property-based testing, deterministic
sweeps, and mutation testing.
The shipped library is axiom-clean against
$\{\texttt{propext}, \texttt{Classical.choice}, \texttt{Quot.sound}\}$
and pinned to Lean 4.28 + Mathlib v4.28.0. The cross-prover ladder
is offline-first and free of large language models (LLMs). SMT and
ATP backends are deterministic decision procedures with verifiable
Lean reconstruction. LLM-augmented closure lives in a separate
companion software development kit (SDK) that the library does not
depend on. Hard structural Mathlib gaps are optionally routed to
the Aristotle automated theorem prover
\citep{harmonic2025aristotle} for development-time restocking only,
with kernel-checked reconstruction inside Lean and an audit gate
against vacuous closures.

\paragraph{Case study (vacuous-truth catch).}
On 2026-04-26 the Aristotle prover closed a candidate theorem in
our \texttt{wald\_wolfowitz\_optimal} development by routing
through a contradictory premise. A Bool partition combined with a
log-boundary constraint forced both $\alpha \geq 1$ and $\alpha <
1$, closing the goal vacuously. The closure was kernel-checked.
The theorem statement was the bug. The library audit caught the
issue and forced a restatement to the two-measure form. We treat
this kind of catch as a first-class library responsibility, not as
a separate review pass.

\paragraph{Paper organisation.}
Section~\ref{sec:tactic-surface} describes the ten-tactic surface
and the registry attributes.
Section~\ref{sec:library} groups the shipped library content by
domain. Section~\ref{sec:aristotle} describes the Aristotle
integration as a development-time oracle and the vacuous-truth
refusal policy.
Section~\ref{sec:discussion} compares against aesop, Sledgehammer,
CoqHammer, and Mathlib's existing tactic suite, states limitations,
and reports a small head-to-head between Aristotle, the DeepSeek
Prover-V2 (DSPv2) model, and Claude Opus 4.6 on a subset of
FormalAVS \citep{anonymous2026formalavs}. Full numbers are in the
dataset paper.

\paragraph{Reproduction and code.}
The \texttt{pythia} library is open source. The URL is redacted
for double-blind review. Axiom audit transcripts and reproduction
recipes are listed in Appendix~\ref{sec:appendix-repro}. The
dataset of theorem-proving benchmarks (FormalAVS,
\citet{anonymous2026formalavs}) is reported in a separate paper.
